Tentu! Berdasarkan kedua sumber dokumen yang diberikan, berikut adalah kumpulan seluruh persamaan matematis metodologi STATIS, yang dituliskan dalam format LaTeX dan dikelompokkan berdasarkan makalahnya:

### 1. Persamaan dari Makalah 2024 ("STATIS methodology: An overview")

Makalah ini berfokus pada pengembangan kemampuan inferensi, sehingga banyak memaparkan pemodelan matriks antar studi dan vektor struktur:

*   **Matriks Data (Tipe 1 dan 2):**
    Untuk studi dengan objek yang sama (Tipe 1):
    $$A_i = X_i D_{pi} X_i^t D_n, \quad i = 1, \dots, k$$
    Untuk studi dengan variabel yang sama (Tipe 2):
    $$B_i = X_i^t D_{ni} X_i D_p, \quad i = 1, \dots, k$$

*   **Matriks Interstruktur (Interstructure Matrix):**
    $$S_{ij} = \text{tr}(A_i A_j^t), \quad i=1,\dots,k, \quad j=1,\dots,k$$
    Produk skalar Hilbert-Schmidt untuk masing-masing tipe:
    $$(A_i|A_j) = \text{tr}(A_i A_j^t), \quad i=1,\dots,k$$
    $$(B_i|B_j) = \text{tr}(B_i B_j^t), \quad i=1,\dots,k$$

*   **Dekomposisi Spektral Matriks Interstruktur ($S$):**
    $$S = \sum_{j=1}^k \lambda_j \alpha_j \alpha_j^t$$

*   **Pemodelan Vektor Struktur (Pengembangan Baru):**
    Persamaan model dasar struktur umum dengan estimasi residu/galat ($E$):
    $$\beta_i \alpha_i^t + E$$
    Di mana vektor struktur $\beta_i$ diestimasi sebagai:
    $$\beta_i = \lambda_i \alpha_i, \quad i = 1, \dots, h$$

---

### 2. Persamaan dari Makalah 1994 ("The ACT (STATIS method)")

Makalah ini berisi pondasi teknis yang mendalam mengenai jarak Euclidean, produk skalar, dan aljabar multivariat.

*   **Pemusatan (Centering) dan Pembobotan Observasi:**
    Syarat bobot matriks diagonal:
    $$\sum m_i = 1$$
    Kendala pemusatan pada matriks:
    $$\sum_i m_i X_{ij} = 0$$
    Matriks kovarians pada tahap $k$:
    $$V_k = X_k' D_k X_k$$

*   **Produk Skalar Antar Observasi dan Variabel:**
    $$W = X X'$$
    $$W = X Q X'$$
    $$(x_i|x_{i'}) = \sum_j X_{ij} X_{i'j}$$
    $$(x_j|x_{j'}) = \sum_i m_i X_{ij} X_{ij'}$$
    $$(x_i|x_{i'}) = \sum_j \sum_{j'} Q_{jj'} X_{ij} X_{i'j'}$$
    $$(e_j|e_{j'}) = Q_{jj'}$$ dan $$||e_i||^2 = m_i$$
    Posisi variabel pada basis ortogonal:
    $$x_j = \sum_i X_{ij} e_i$$

*   **Jarak Euclidean (Euclidean Distances):**
    Hubungan jarak dan produk skalar:
    $$d^2(i, i') = \sum_j (X_{ij} - X_{i'j})^2$$
    $$d^2(i, i') = ||x_i - x_{i'}||^2$$
    $$W_{ii'} = \overrightarrow{OM_i} \cdot \overrightarrow{OM_{i'}}$$
    $$||\overrightarrow{M_i M_{i'}}|| = d(i, i')$$
    $$d^2(i, i') = W_{ii} + W_{i'i'} - 2W_{ii'}$$
    $$W_{ii'} = \frac{1}{2} [d^2(0, i) + d^2(0, i') - d^2(i, i')]$$

*   **Jarak Antar Matriks Skalar Produk ($W_k$):**
    $$||W_1 - W_2||^2 = \sum_i \sum_{i'} [(W_1 - W_2)_{ii'}]^2 = \text{Tr}[(W_1 - W_2)^2]$$
    Versi terbobot:
    $$||W_1 - W_2||^2 = \sum_i \sum_{i'} S_{ii} S_{i'i'} [(W_1 - W_2)_{ii'}]^2$$
    Rumus umum jarak matriks:
    $$||W_1 - W_2||^2 = \text{Tr}[(W_1 S - W_2 S)^2]$$
    Jarak interstruktur $W_k$ ke $W_{k'}$:
    $$d^2(W_k, W_{k'}) = \sum_i \sum_{i'} [(W_k - W_{k'})_{ii'}]^2$$
    $$d^2(W_k, W_{k'}) = \sum_i \sum_{i'} S_{ii} S_{i'i'} [(W_k - W_{k'})_{ii'}]^2$$
    $$d^2(W_k, W_{k'}) = \text{Tr}[(W_k S - W_{k'} S)^2]$$

*   **Koefisien RV & Skalar Produk Matriks:**
    Bentuk murni dan normanya:
    $$\text{Tr}(W_k D W_{k'} D)$$ atau $$\frac{\text{Tr}(W_k D W_{k'} D)}{\sqrt{\text{Tr}(W_k D)^2} \sqrt{\text{Tr}(W_{k'} D)^2}}$$
    $$(W_k|W_{k'}) = \text{Tr}(W_k W_{k'}) = \text{Tr}(W_k S W_{k'} S)$$
    Versi dinormalkan:
    $$\frac{W_k}{||W_k||} = \frac{W_k}{\sqrt{\text{Tr}(W_k S)^2}}$$
    $$\left( \frac{W_k}{||W_k||} \Big| \frac{W_{k'}}{||W_{k'}||} \right) = \frac{\text{Tr}(W_k S W_{k'} S)}{\sqrt{\text{Tr}(W_k S)^2} \sqrt{\text{Tr}(W_{k'} S)^2}}$$

*   **Matriks Kompromi & Dekomposisi:**
    Rumus dasar Matriks Kompromi:
    $$W = \sum_{k=1}^K \alpha_k W_k$$
    Optimasi kuosien Rayleigh:
    $$\frac{x' W W^* x}{||x||^2}$$
    Persamaan residu (*loss function*):
    $$||W - W_h||^2 = \sum_k \sum_{k'} [(W - W_h)_{kk'}]^2 = \sum_{l=h+1}^r \sigma_l^2$$
    $$||W - W_h||^2 = \text{Tr}[(W S - W_h S)^2] = \sum_{l=h+1}^r \sigma_l^2$$
    Dekomposisi matriks kompromi (SVD):
    $$W = \sigma_1 p_1 p_1' + \dots + \sigma_r p_r p_r' = P \Sigma P'$$
    Syarat ortogonal:
    $$P'P = I$$
    Hubungan sudut koordinat:
    $$\overrightarrow{OM_k} \cdot \overrightarrow{OM_{k'}} = (W_h)_{kk'}$$

*   **Lampiran: Properti Ruang Vektor & Operator Adjoin:**
    $$(\mathscr{S}(x))(y) = (x|y)_s$$
    $$(u|v)_{s^*} = v(\mathscr{S}^{-1}(u))$$
    $$(\mathscr{A}'(v))(x) = v(\mathscr{A}(x))$$
    $$(x|\mathscr{A}^*(y))_{s_J} = (\mathscr{A}(x)|y)_{s_I}$$
    $$\mathscr{A}^* = \mathscr{S}_J^{-1} \mathscr{A}' \mathscr{S}_I$$
    Kondisi definit positif:
    $$(\mathscr{A}^*\mathscr{A}(x)|x)_{s_J} \geq 0$$
    Singular Value Decomposition (SVD) pemetaan linier:
    $$A = \sigma_1 p_1 q_1' S_J + \dots + \sigma_r p_r q_r' S_J$$
    $$A = P \Sigma Q' S_J$$ dengan $$P' S_I P = I$$ dan $$Q' S_J Q = I$$
    Jarak antar pemetaan linier:
    $$(\mathscr{A}|\mathscr{B}) = \text{Tr}(\mathscr{A}\mathscr{B}^*) = \text{Tr}(\mathscr{A}\mathscr{S}_I^{-1}\mathscr{B}'\mathscr{S}_I)$$
    $$\text{d}^2(\mathscr{A}, \mathscr{B}) = ||\mathscr{A} - \mathscr{B}||^2 = \text{Tr}((\mathscr{A}-\mathscr{B})\mathscr{S}_I^{-1}(\mathscr{A}-\mathscr{B})'\mathscr{S}_I)$$
    $$\min_{\mathscr{B} \in L(E,F), \text{rank of } \mathscr{B} \leq h} ||\mathscr{A} - \mathscr{B}||^2 = ||\mathscr{A} - \mathscr{A}_h||^2 = \sum_{l=h+1}^r \sigma_l^2$$